%% fig-dataflow.tex -- Figure: data flow and state ownership.
%%
%% Three stores, three owners, and the two reconciliation boundaries that exist
%% because no transaction can span them. Solid edges are transactional; the two
%% dashed edges are the crossings that can diverge and must be reconciled.

\begin{figure}[tbp]
\centering
\ifcbexcalidraw
  \begin{cbwide}\cbdiagram{fig-dataflow}\end{cbwide}
\else
\begin{cbwide}
\begin{tikzpicture}[x=1cm,y=1cm]

  % ======================= writers ==========================================
  \node[anchor=north west, font=\figB\bfseries, text=cbMuted] at (0,0.44) {WRITERS};
  \node[surface, anchor=north west, text width=3.34cm, minimum height=0.80cm]
    (rh) at (0,0.10)
    {\textbf{Route handler}\\{\figB inline validation, then\\durable mutation}};
  \node[async, anchor=north west, text width=3.34cm, minimum height=0.80cm]
    (wk) at (0,-1.06)
    {\textbf{Background loop}\\{\figB a claimed job, writing\\the same tables}};
  \node[control, anchor=north west, text width=3.34cm, minimum height=0.80cm]
    (apl) at (0,-2.22)
    {\textbf{Apply + promoter}\\{\figB the release path}};

  % ======================= the three owners =================================
  \node[anchor=north west, font=\figB\bfseries, text=cbMuted] at (4.64,0.54)
    {AUTHORITATIVE OWNERS};

  \node[db, anchor=north west, text width=3.86cm, minimum height=3.02cm,
        align=left] (rel) at (4.64,0.10)
    {\textbf{Relational metadata}\\[0.35ex]
     {\figB \textsc{owns} the control plane ---\\
      registries \dt jobs \dt approvals\\
      runs \dt deployments \dt verdicts\\
      audit log \dt effect ledger\\
      and the file \emph{inventory}\\[0.5ex]
      \textsc{not} file bytes\\
      \textsc{not} prompt versions}};

  \node[db, anchor=north west, text width=3.94cm, minimum height=1.32cm,
        align=left] (obj) at (9.32,0.10)
    {\textbf{Object storage}\\[0.2ex]
     {\figB \textsc{owns} file \emph{bytes} --- uploads,\\
      artifacts, exports, archived traces\\[0.4ex]
      \textsc{not} the inventory}};

  \node[extern, anchor=north west, text width=3.94cm, minimum height=1.42cm,
        align=left] (mlf) at (9.32,-1.84)
    {\textbf{\mlflow{}}\\[0.2ex]
     {\figB \textsc{owns} prompt registry versions,\\
      aliases, traces, assessments\\[0.3ex]
      \textsc{not} workflow or tool metadata}};

  % ======================= relational domains ===============================
  \node[anchor=north west, font=\figB\bfseries, text=cbMuted] at (13.94,0.54)
    {RELATIONAL DOMAINS};
  \foreach [count=\i from 0] \d/\t in {%
      {Core governance}/{items, jobs, approvals},
      {Registries}/{prompt, tool, skill, MCP},
      {Knowledge \& quality}/{bases, chunks, judges},
      {Platform}/{projects, files, secrets}}
  {
    \node[muted, anchor=north west, text width=3.52cm, minimum height=0.74cm,
          align=left, font=\figB] at (13.94, {0.10-\i*0.80})
      {\textbf{\d}\\{\figB\t}};
  }

  % ======================= edges ============================================
  \draw[flow] (rh.east)  -- (rel.west);
  \draw[flow] (wk.east)  -- (rel.west);
  \draw[flow] (apl.east) -- (rel.west);
  \node[anchor=south, font=\figB, text=cbInk, inner sep=1.5pt] at (4.20,-0.90)
    {\rotatebox{90}{\makebox[0pt][c]{same tables}}};

  % The two crossings that no transaction can cover. The double bar at midpoint is
  % the conventional mark for a boundary rather than a link, so the edge says what
  % it is without depending on the caption.
  \draw[boundary] (rel.east) -- (obj.west);
  \draw[boundary] (rel.east) -- (mlf.west);

  % The right-hand column is a partition of the relational store, not a fourth
  % store. A connector back to it would have to cross the MLflow box, so the
  % relationship is carried by the column header and the caption instead.
  \draw[weak, draw=cbMuted] (13.62,0.10) -- (13.62,-3.02);

  % ======================= annotations ======================================
  \node[govern, anchor=north west, text width=8.60cm, align=left,
        minimum height=0.94cm] at (4.64,-3.42)
    {\textbf{The two dashed crossings are dual-write boundaries}\\[0.2ex]
     {\figB A database-resident mutation and its \code{audit\_record} call share
      \emph{one} SQL transaction, so control-plane state and its audit row cannot
      diverge. Object-storage writes and external registry or provider effects
      cannot join that transaction: they are settled by idempotent retry and
      reconciliation, not by atomicity (\secref{sec:design}). Importing the audit
      helper into a module does not prove every mutation in it is audited.}};

  \node[store, anchor=north west, text width=3.52cm, align=left,
        minimum height=0.94cm] at (13.94,-3.42)
    {\textbf{Reader path}\\[0.2ex]
     {\figB A read of a \gasset resolves relational metadata first, then
      dereferences bytes or registry versions. The inventory is therefore never
      rebuilt from a bucket listing.}};

\end{tikzpicture}
\end{cbwide}
\fi
\caption{Data flow and state ownership. Route handlers, background loops, and the
release path all write the same relational tables --- which are authoritative for
the control plane \emph{and} for the file inventory. Object storage owns file
bytes and nothing else; \mlflow{} owns prompt registry versions, aliases, traces,
and assessments. Solid edges are transactional. The two dashed edges are the
architecture's honest limit: no transaction spans \sys's database and an external
store, so those crossings are reconciled rather than made atomic. Naming them
explicitly is what lets the audit guarantee in \secref{sec:components} be stated
precisely instead of gestured at --- database-resident mutations and their audit
rows share one transaction, and nothing wider is claimed. The right-hand column
is a partition of the relational store, not a fourth store.}
\label{fig:dataflow}
\figkey
\end{figure}
